\documentclass[12pt]{article}
\usepackage[T1]{fontenc}
\usepackage[utf8]{inputenc}
\usepackage{lmodern}
\usepackage{textcomp}
\usepackage{enumitem}
\usepackage{geometry}
\usepackage{graphicx}
\usepackage{amsmath, amssymb}
\geometry{margin=1in}
\begin{document}
\begin{center}
Political Science - Graduate Level Assessment 7 \
Total Marks: 40 \
Answer all questions.
\end{center}

\section*{Multiple Choice (10 marks)}
\begin{enumerate}[label=\arabic*.]
\item What was the 'New Populism'?
\begin{enumerate}[label=(\alph*)]
    \item A strand of neo-isolationist sentiment
    \item A strand of internationalist sentiment
    \item An expression of American cultural superiority
    \item Increased incorporation of public opinion in foreign policy making
\end{enumerate}
\item Which of the following are trends to emerge in the global defence trade since the end of the Cold War?
\begin{enumerate}[label=(\alph*)]
    \item A growing effort by governments to prevent the spread of weapons, especially of WMDs.
    \item A growth in volume of the arms trade.
    \item A change in the nature of the defence trade, linked to the changing nature of conflict.
    \item All of these options.
\end{enumerate}
\item Historical materialism is founded on the ideas and philosophies of which theorists?
\begin{enumerate}[label=(\alph*)]
    \item Barry Buzan and Ole Waever
    \item Kenneth Waltz and Hans Morgenthau
    \item Karl Marx and Friedrich Engels
    \item Adam Smith and Karl Marx
\end{enumerate}
\item The departments of the executive branch that assist the president in designing and carrying out U.S. foreign policy are known as
\begin{enumerate}[label=(\alph*)]
    \item the United Nations.
    \item the National Security Council.
    \item the State Department.
    \item the National Security Agency.
\end{enumerate}
\item Why is there so much uncertainty over which states have nuclear weapons?
\begin{enumerate}[label=(\alph*)]
    \item Leaders have incentives to lie
    \item If leaders revealed their programs, they would be more likely to be attacked
    \item Leaders will not always grant foreign monitors access to their nuclear programs
    \item ALL of the above
\end{enumerate}
\end{enumerate}

\section*{True / False (10 marks)}
\textit{Answer True or False}
\begin{enumerate}[label=\arabic*.]
\setcounter{enumi}{5}
\item True or False: China is considered a state actor and does not fit the classification of nonstate actors that threaten the United States.
\item True or False: American exceptionalism encourages the use of tariffs to protect domestic industries from international competition.
\item True or False: Under George H.W. Bush and Bill Clinton, Congress did not assert its primacy in foreign policy, reflecting a different power dynamic.
\item True or False: Ideational approaches to US foreign policy during the Cold War emphasize ideology and beliefs more than material interests.
\item True or False: Wealth and GDP are the most significant predictors of nuclear weapons acquisition by states.
\end{enumerate}

\section*{Long Form Response (20 marks)}
\textit{Answer each question with a clear and complete explanation.}
\begin{enumerate}[label=\arabic*.]
\setcounter{enumi}{10}
\item Discuss the broader political considerations that led to the institutional innovation of security studies, focusing on the shift from tactical operations to long-term strategic planning.
\item Discuss the concept of the 'subaltern' in postcolonialism, analyzing its implications for marginalized populations within hegemonic power structures.
\end{enumerate}

\vfill
\noindent\textit{End of Paper}
\end{document}